\documentclass[a4paper,11pt,openany]{book}
\usepackage[utf8]{inputenc}
\usepackage[left=2.54cm,top=2.54cm,right=2.54cm,bottom=2.54cm]{geometry}
\usepackage[spanish]{babel}
\usepackage{amsmath}
\usepackage{amsfonts}
\usepackage{amssymb}
\usepackage{graphicx}
\usepackage{color}
\usepackage[usenames,dvipsnames]{xcolor}
\usepackage{pifont}
\usepackage{marvosym}
\newtheorem{teo}{Teorema}
\newtheorem{ejemplo}{Ejemplo}
\newtheorem{defi}{Definición}
\newtheorem{coro}{Corolario}
\newtheorem{prueba}{Prueba}
\newtheorem{exmp}{Example}[section]
\newtheorem{ejer}{Ejercicio}[section]
\def\proof{\paragraph{\textsf{Demostración.} }}
\def\endproof{\hfill $\blacksquare$ \\}
\usepackage{multirow, array} % para las tablas
\usepackage{multirow}
\usepackage{tabularx}
\usepackage{float} % para usar [H]
\usepackage{tikz}
\usepackage[all]{xy}
\usepackage{cancel}
\usetikzlibrary{positioning}
\usepackage{enumitem}
\newcommand*{\itembolasazules}[1]{% bolas 3D
\footnotesize\protect\tikz[baseline=-3pt]%
\protect\node[scale=.7, circle, shade, ball
color=green]{\color{white}\Large\bf#1};}
\usepackage{tcolorbox} 
\tcbuselibrary{listingsutf8}
\newtcolorbox[auto counter,number within=section]{example}[2][]
{colback=green!5!white,colframe=green!75!black,fonttitle=\bfseries, title=Ejercicio~\thetcbcounter: #2,#1}
\usepackage{background}
\backgroundsetup{
placement=center,
angle=0,
scale=1.1,
contents= {{\includegraphics{HojaCuadriculada.png}}}
}
 
 
\begin{document}
\begin{titlepage}
 
\begin{center}
\vspace*{-1in}
\begin{figure}[htb]
\begin{center}
\includegraphics[width=7cm]{ETITC.png}
\end{center}
\end{figure}
 
 
{\sc \huge Escuela Tecnológica Instituto Técnico Central (ETITC)}\\
\vspace*{0.15in}
Facultad de sistemas\\
\vspace*{0.6in}
\begin{Large}
\textbf{Taller 3: Conjuntos en el Plano Complejo} \\
\textbf{Matem{\'a}ticas Especiales}\\
\end{Large}
\vspace*{0.3in}
\begin{large}
{\bf Autores} \\
 
\ 
 
Sergio Alejandro Enrrique Caballero Leon\\ 
Johan Alejandro Sogamoso Camacho \\
David Andrés Valero Vanegas \\
\end{large}
\vspace*{0.3in}
 
\end{center}
 
\begin{center}
{\bf Presentado a:} \\
 
\ 
 
Carlos Romero \\
 
\
 
Bogot{\'a}, Septiembre de 2022.
\end{center}
 
\end{titlepage}

\newpage

\definecolor{ao(english)}{rgb}{0.0, 0.5, 0.0}

\graphicspath{ {images/} }

\begin{center}
\textbf{Conjuntos}
\end{center}

En los ejercicios \textbf{(\,1\,)} al \textbf{(\,12\,) graficar}, cada uno en un plano complejo, el conjunto de puntos $\mathnormal{S}$ que satisface la igualdad o desigualdad dada y \textbf{escribir qué tipo de conjunto es:} Círculo, Disco, Anillo, etc. Además decir si el conjunto es Abierto, Cerrado ó Semi-Abierto (ó Semi-Cerrado).\\

\textcolor{ao(english)}{(\,1\,)} $\bf{|z\,-\,4\,+\,3\,i|\,=\,5}$.

\textcolor{ao(english)}{(\,2\,)} $\bf{|z\,+\,2\,-\,5\,i|\,\leq\,4}$.

\textcolor{ao(english)}{(\,3\,)} $\bf{2\,\leq\,|z\,+\,3\,i|\,<\,6}$.

\textcolor{ao(english)}{(\,4\,)} $\bf{2\,<\,|z\,+\,5\,-\,2\,i|\,<\,9}$.

\textcolor{ao(english)}{(\,5\,)} $\bf{|z\,-\,4\,+\,8\,i|\,<\,\dfrac{7}{2}}$.

\textcolor{ao(english)}{(\,6\,)} $\bf{|z\,+\,4\,+\,2\,i|\,=\,6}$.

\textcolor{ao(english)}{(\,7\,)} $\bf{0\,<\,|z\,-\,1\,+\,4\,i|\,\leq\,5}$.

\textcolor{ao(english)}{(\,8\,)} $\bf{2\,\leq\,|z\,-\,1|\,\leq\,9}$.

\textcolor{ao(english)}{(\,9\,)} $\bf{\mathnormal{Im}\,(\overline{z}\,+\,3\,i)\,=\,6}$.

\textcolor{ao(english)}{(\,10\,)} $\bf{\mathnormal{Re}\,(z\,+\,8)\,=\,5}$.

\textcolor{ao(english)}{(\,11\,)} $\bf{\mathnormal{Im}\,(z\,-\,10\,+\,2\,i)\,>\,\dfrac{5}{4}}$.

\textcolor{ao(english)}{(\,12\,)} $\bf{\mathnormal{Re}\,(1\,+\,7\,i\,+\,z)\,>\,\dfrac{2}{3}}$.

\end{document}